\section{Functional implementation}

\subsection{Floating-point arithmetic subsystem}

The arithmetic subsystem is the computational backbone of the solver. The design philosophy is to build one reusable IEEE-754 FP32 arithmetic library, then let the top-level controller schedule these blocks repeatedly to evaluate the cubic-solving algorithm. In other words, the project follows a \textbf{shared datapath + control FSM} strategy instead of instantiating a separate arithmetic network for each formula. This approach has three advantages:
\begin{itemize}[noitemsep]
    \item it keeps the architecture structurally regular and easier to debug
    \item it allows the same arithmetic blocks to be reused in coefficient normalization, discriminant computation, iterative approximation, and final root reconstruction
    \item it makes the behavior of the hardware closely match the golden model, where each high-level formula is also decomposed into a sequence of primitive operations
\end{itemize}

\subsubsection{Adder module}
The \texttt{Add\_FP} module performs both addition and subtraction on FP32 operands.
\begin{itemize}[noitemsep]
    \item \textbf{Operand alignment}: the exponents are compared first so that the smaller mantissa can be right-shifted before accumulation.
    \item \textbf{Unified add/sub datapath}: subtraction is implemented implicitly by changing operand signs and by using the internal \texttt{op} control signal.
    \item \textbf{Normalization and rounding}: after mantissa accumulation, the result is normalized and packed back into IEEE-754 format.
\end{itemize}

\subsubsection{Multiplier module}
The \texttt{Mul\_FP} module computes the product of two FP32 numbers through a pipelined mantissa multiplier followed by normalization and rounding.
\begin{itemize}[noitemsep]
    \item \textbf{Sign and exponent processing}: the output sign is obtained by XOR, while the output exponent is formed by exponent addition and bias correction.
    \item \textbf{Mantissa product generation}: the 24-bit mantissas are multiplied with a staged partial-product architecture.
    \item \textbf{High reuse}: the module is repeatedly used to evaluate powers, cross-products, and complex arithmetic terms during Pad\'e iterations.
\end{itemize}

\subsubsection{Divider module}
The \texttt{Div\_FP} module evaluates one FP32 division through a radix-4 pipelined divider.
\begin{itemize}[noitemsep]
    \item \textbf{Pre-processing}: sign, zero cases, division-by-zero, NaN, and exponent difference are identified before quotient generation.
    \item \textbf{Radix-4 quotient generation}: each stage emits two quotient bits, which provides a predictable latency suitable for controller scheduling.
    \item \textbf{Round-pack stage}: the quotient is converted back into a normalized IEEE-754 result before being registered at the output.
\end{itemize}

\subsection{Main solver behavior}

The \texttt{Cubic\_Solver} module is a controller-centric architecture. Its main responsibility is not to compute every formula in parallel, but to orchestrate the arithmetic modules so that the correct operands are presented to \texttt{Add\_FP}, \texttt{Mul\_FP}, and \texttt{Div\_FP} at the correct cycle offsets. Each state of the FSM corresponds to one algorithmic phase, while \texttt{delay\_count} acts as a micro-scheduler that determines which arithmetic expression is being issued in that phase.

Functionally, the top-level solver performs the following sequence:
\begin{enumerate}
    \item Capture input coefficients and standardize them by dividing with $-3a$.
    \item Compute $\Delta_0$, $\Delta_1$, and $\Delta$.
    \item Classify the equation into one of the three solution branches.
    \item Execute the square-root iteration for $\sqrt{\Delta}$ or $\sqrt{-\Delta}$.
    \item Execute the cube-root iteration in either real or complex mode to obtain $C$.
    \item Combine $\hat{b}$, $C$, $\overline{C}$, and $\Delta_0/C$ to generate the three final roots.
\end{enumerate}

At the implementation level, this sequence is realized by repeatedly reconfiguring the inputs of the three arithmetic blocks:
\begin{itemize}[noitemsep]
    \item \texttt{MulInput1/2} carry products such as $\hat{b}^2$, $\hat{b}^3$, $\Delta_1^2$, $P^2$, $P^3$, and the real/imaginary cross-products used in complex cube-root iteration.
    \item \texttt{AddInput1/2} combine intermediate terms into expressions such as $\Delta_0$, $\Delta_1$, $3S$, $3P^2+S$, $2P^3+S$, and the final root formulas.
    \item \texttt{DivInput1/2} are used whenever the algorithm requires normalization by division, Pad\'e update evaluation, or the quotient $\Delta_0/C$.
\end{itemize}

This means the top-level block is best interpreted as a scheduled arithmetic engine rather than a static combinational solver.

\subsection{Main control signals}

The most important control intent of the design is concentrated in a small set of scheduling and branch signals. Their functional roles are summarized in Table~\ref{tab:control_signals}.

\begin{table}[H]
\centering
\caption{Main control and scheduling signals in \texttt{Cubic\_Solver}}
\label{tab:control_signals}
\renewcommand{\arraystretch}{1.2}
\begin{tabular}{p{3.0cm} p{2.0cm} p{1.4cm} p{8.0cm}}
\toprule
\textbf{Signal} & \textbf{Type} & \textbf{Width} & \textbf{Functional meaning} \\
\midrule
\texttt{state} & Reg & 4 & Current state of the main finite state machine. Selects the active computational phase. \\
\texttt{next\_state} & Reg & 4 & Next-state combinational result derived from the current state, counters, and branch conditions. \\
\texttt{case\_index} & Reg & 2 & Encodes the selected solution branch: negative discriminant, simplified real branch, or general real branch. \\
\texttt{delay\_count} & Reg & 7 & Fine-grain schedule counter used to align operand issue with the latency of arithmetic pipelines. \\
\texttt{iteration\_count} & Reg & 5 & Counts the number of Pad\'e iterations already completed for square-root or cube-root evaluation. \\
\texttt{READY\_BRANCH\_CHECK} & Reg & 1 & Internal synchronization flag indicating that the branch-specific cube-root input and initial estimate are ready. \\
\texttt{MulInput1}, \texttt{MulInput2} & Reg & 32 each & Operands dispatched to the shared floating-point multiplier at the current cycle. \\
\texttt{AddInput1}, \texttt{AddInput2} & Reg & 32 each & Operands dispatched to the shared floating-point adder at the current cycle. \\
\texttt{DivInput1}, \texttt{DivInput2} & Reg & 32 each & Operands dispatched to the shared floating-point divider at the current cycle. \\
\texttt{cbrt\_scale\_in} & Reg & 8 & Exponent-related helper input used to construct the initial estimate for cube-root iteration. \\
\texttt{exponent\_max} & Reg & 8 & Stores the larger exponent between the real and imaginary parts when a complex cube-root seed is required. \\
\texttt{done} & Reg & 1 & Output completion flag asserted in the \texttt{DONE} state. \\
\bottomrule
\end{tabular}
\end{table}

\subsection{The internal registers}

Besides the explicit control signals, the solver uses a compact set of internal registers to hold operands and intermediate results between states:
\begin{itemize}[noitemsep]
    \item \texttt{r\_a}, \texttt{r\_b}, \texttt{r\_c}, \texttt{r\_d}: working coefficients. After the \texttt{STANDARDIZE} state, these registers no longer contain the original input values but the normalized coefficients used by the algorithm.
    \item \texttt{r\_delta\_0}, \texttt{r\_delta\_1}, \texttt{r\_delta}: algebraic invariants used for branch classification and final-root reconstruction.
    \item \texttt{r\_sqrt\_input}, \texttt{r\_sqrt\_output}: operand and current estimate for Pad\'e square-root iteration.
    \item \texttt{r\_cbrt\_re\_input}, \texttt{r\_cbrt\_im\_input}: real and imaginary parts of the cube-root operand.
    \item \texttt{r\_cbrt\_re\_output}, \texttt{r\_cbrt\_im\_output}: current cube-root estimate.
    \item \texttt{r\_C\_re}, \texttt{r\_C\_im}: final value of $C$ after the cube-root stage completes.
    \item \texttt{r\_temp[0:63]}: scratchpad register bank used to preserve pipeline results across multiple delay slots in long arithmetic schedules.
\end{itemize}

\subsection{Functional role of branch control and initialization}

The signal \texttt{case\_index} has a central functional meaning because it determines which numerical path the hardware will follow after the discriminant values have been computed:
\begin{itemize}[noitemsep]
    \item \texttt{2'b01}: discriminant is negative, so the solver must compute a complex-valued cube root and later generate roots from $C$ and $\overline{C}$.
    \item \texttt{2'b10}: special real branch with simplified cube-root evaluation.
    \item \texttt{2'b11}: general real branch, where the solver computes both $C$ and $\Delta_0/C$ before forming the three roots.
\end{itemize}
Because all later states depend on this 2-bit decision, the \texttt{CASE\_CHECK} and \texttt{BRANCH\_CHECK} states serve as the bridge between the algebraic invariants and the branch-specific datapath schedule.

Another important design feature is the use of \textbf{hardware-friendly initial estimates} for the iterative square-root and cube-root blocks.

For the square-root stage, the initial approximation is formed directly from the exponent field of the operand. The wire
\begin{equation}
\texttt{sqrt\_output\_exp = ((exp(\Delta) + 127) >> 1)}
\end{equation}
approximates the exponent of $\sqrt{\Delta}$, and the corresponding mantissa is copied to create \texttt{sqrt\_output\_init}. This gives the Pad\'e square-root iteration a starting point that is already close in magnitude to the true solution.

For the cube-root stage, the design uses the helper block \texttt{Divide\_int} to approximate the exponent transform
\begin{equation}
E_{\sqrt[3]{S}} \approx 127 + \frac{E_S - 127}{3}
\end{equation}
without invoking an additional floating-point divider. In the complex branch, the design first compares the exponents of the real and imaginary parts, selects the larger one as \texttt{exponent\_max}, and uses it to build a conservative initial estimate for the complex cube root. In the real branches, the exponent of the real operand alone is sufficient.

This initialization strategy is important because it reduces the number of Pad\'e iterations needed for convergence and makes the iterative states more stable in hardware.
